\anexo{Transcripción de la entrevista a Sebastián López}
\begin{description}[leftmargin=0pt, itemindent=\parindent, labelsep=0.5em, itemsep=0pt, parsep=0pt, topsep=0pt]
    \item[\textbf{Entrevistador:}] Bueno, antes de arrancar, te agradezco Sebastián por tomarte el tiempo. Te cuento brevemente de qué trata el proyecto. Estoy desarrollando una propuesta de sistema para observar movimientos del tren inferior en jugadores de fútbol, principalmente rodilla, tobillo y apoyos. La idea no es diagnosticar lesiones ni reemplazar al cuerpo médico, sino analizar videos de entrenamientos y competiciones para detectar posibles patrones de riesgo y generar información que pueda servir como apoyo para la prevención. Para comenzar, ¿te podrías presentar brevemente?

    \item[\textbf{Sebastián López:}] Sí, dale. Hola, soy Sebastián López, arquero de fútbol profesional del Club Atlético Los Andes. Estoy acá para brindar un poco de mi experiencia en cuanto a entrenamientos, partidos, funcionalidad y todo lo relacionado con los mecanismos de movimiento.

    \item[\textbf{Entrevistador:}] Gracias por la presentación. ¿Cuál es tu edad y hace cuánto jugás al fútbol de manera competitiva?

    \item[\textbf{Sebastián López:}] Bueno, tengo 40 años y de manera competitiva, desde los 14 años, cuando vine a Buenos Aires a fichar en las divisiones inferiores de Banfield. Desde ahí arranqué y no paré más.

    \item[\textbf{Entrevistador:}] Claro. Estuve viendo que tuviste bastante recorrido de clubes en tu carrera.

    \item[\textbf{Sebastián López:}] Sí. Arranqué acá en Argentina. Jugué en dos clubes: Banfield, en mis inicios, y después Los Andes. Ahora ya es mi cuarto año en el club. También estuve 13 años en el exterior, jugando en Chile y Colombia. Después decidí volver, hace cuatro años, y me abrió las puertas el Club Los Andes. Acá estamos.

    \item[\textbf{Entrevistador:}] Muy interesante. ¿Siempre jugaste de arquero?

    \item[\textbf{Sebastián López:}] Sí, siempre. Desde muy chico arranqué baby fútbol en San Nicolás, con 6 años, en el Club Don Bosco. Después, a los 11 años, pasé a jugar en cancha de 11. A los 13 empecé a probarme en clubes de primera división, como Newell's, que me quedaba cerca de mi ciudad. Después vine a Buenos Aires un par de veces y, a los 14 años, vine definitivamente a Banfield. Me quedé en la pensión y ahí empezaron los cambios en cuanto a la preparación física. Empezás a entrenar de otra manera, prácticamente todos los días, a incrementar la masa muscular mediante gimnasio y con entrenamientos mucho más intensos. Entonces, vas adaptando el cuerpo para todo lo que viene.

    \item[\textbf{Entrevistador:}] Entiendo. Entonces, por lo que me comentás, el entrenamiento semanal incluye gimnasio todos los días. ¿Qué otras actividades incluye?

    \item[\textbf{Sebastián López:}] Sí. Ahora, por ahí, ya no se hace tanta carga de peso, sino movimientos mucho más específicos. Por ejemplo, movilidad de tobillos, movilidad de rodillas, muchos trabajos de estabilidad con bandas y muchos trabajos de propiocepción. Más que nada, se busca simular o acercarse lo más posible a situaciones reales de competencia. Por ejemplo, saltás a buscar una pelota, te mueven en el aire y tenés que caer. Con la propiocepción entrenada durante la semana, el cuerpo, la rodilla y el tobillo actúan de una manera diferente, más fortalecida. No es algo que pase desapercibido. Es un trabajo muy fino, durante la semana y durante mucho tiempo. Por ahí no son trabajos de carga, pero sí de arranques, giros, frenos, saltos a un pie, cambios de dirección. Todo está adaptado para que el cerebro esté preparado, porque en el partido uno hace las cosas sin pensar; ya están prácticamente automatizadas.

    \item[\textbf{Entrevistador:}] Gracias por la explicación. Entonces, por lo que comentás, es como que todo el tiempo trabajan la prevención. Si bien no directamente, sí se trabaja el cuidado físico y la prevención de lesiones.

    \item[\textbf{Sebastián López:}] Sí, prácticamente. Cuando llegás al fútbol profesional, y en mi caso estoy en la parte final de mi carrera, si bien tengo que hacer fuerza para mantener la masa muscular y el tono muscular, se hace mucho más énfasis en la prevención y el fortalecimiento. No se busca tanto ganar masa muscular, sino más que nada mantener y lograr que las articulaciones estén fuertes.

    \item[\textbf{Entrevistador:}] Perfecto. Dentro del Club Los Andes, ¿ustedes usan alguna herramienta para medir el rendimiento o la carga física?

    \item[\textbf{Sebastián López:}] Sí. En mi caso no, porque los arqueros no lo utilizamos, pero los jugadores de campo usan GPS.

    \item[\textbf{Entrevistador:}] ¿Y eso ayuda a saber los niveles de carga durante la semana?

    \item[\textbf{Sebastián López:}] Sí. Lo que marca es una distancia media, una cantidad determinada de arranques, frenos y saltos. Cargado todas las semanas, eso va dando un indicio. No te dice directamente ``hoy corrió de más'', sino que la comparativa la hace el preparador físico día a día. Por ejemplo, la carga del GPS del día de competencia es muy alta en cuanto a frenos y arranques, y suele ser la carga más alta de toda la semana. Durante la semana se trabaja aproximadamente al 70\% u 80\%, y después se bajan las cargas unas 48 horas antes del partido. Sirve para ir midiendo y saber qué jugador está al límite. Muchas veces el jugador te dice que está bien, pero los registros del GPS indican que bajó un poco la intensidad. Entonces ahí se acerca el profe y le pregunta si tiene alguna molestia o fatiga, porque se notó que el registro del GPS dio menos intensidad, menos frenos o menos arranques. Esa es una de las maneras que tienen los profes. También se usa la carga subjetiva. Todos los días, antes y después del entrenamiento, los profes mandan una escala donde el jugador tiene que indicar cómo encontró el entrenamiento: si fue pesado, liviano, de carga media o de carga alta. Así van nivelando y llevando registros de todos los jugadores de manera individual, no tan generalizada. Eso permite que cada jugador tenga un seguimiento y también ayuda a prevenir lesiones.

    \item[\textbf{Entrevistador:}] Gracias por el detalle. En tu caso particular, ¿tuviste alguna lesión de rodilla, tobillo o alguna zona de la pierna?

    \item[\textbf{Sebastián López:}] Sí. De rodilla tuve una lesión hace cuatro años: rotura de meniscos y distensión del ligamento colateral. Tuve que ser intervenido; me retiraron la porción del menisco que estaba dañada. Pero, bueno, a las cinco semanas ya estaba jugando de nuevo. Después tuve distensiones y una rotura total del ligamento colateral interno. Salvo la lesión de menisco, ninguna fue para operar. El ligamento colateral interno se fue regenerando solo, y a las ocho o diez semanas ya estaba jugando nuevamente.

    \item[\textbf{Entrevistador:}] Entiendo. ¿Y en esas lesiones te acordás cómo fue el momento? ¿Fue por contacto o fue solo?

    \item[\textbf{Sebastián López:}] Fue solo. Se me trabó el botín en el pasto, quise girar y se me quedó trabada la rodilla. Por suerte no fue mayor; fue solo el menisco.

    \item[\textbf{Entrevistador:}] ¿Y antes habías sentido alguna molestia, cansancio o alguna señal previa? ¿O fue directamente en el momento por trabarse?

    \item[\textbf{Sebastián López:}] No, fue algo fortuito. Fue una jugada de partido, en un amistoso.

    \item[\textbf{Entrevistador:}] En relación con los arqueros, ¿hay algún preparador físico o persona encargada de observar y corregir aspectos como giros, caídas o frenadas?

    \item[\textbf{Sebastián López:}] Sí, obviamente. El entrenador de arqueros está pendiente de optimizar movimientos, más que nada de manera preventiva. Por ejemplo, hay un día a la semana en el que se coordina la parte de fuerza y se traslada a trabajos de velocidad y reacción. La idea es que el cuerpo vaya adaptando esa fuerza que venimos trabajando. Hacemos trabajos de salto y propiocepción sobre cajones, con desplazamientos, frenos y cambios de dirección. También se adapta al arco: tirarse, ir a buscar una pelota o trabajar el juego aéreo. Pero como estamos en competencia, por ahí son pocos los estímulos de fuerza. Está todo más orientado a la velocidad de reacción y la potencia, no tanto a la carga.

    \item[\textbf{Entrevistador:}] Gracias, se entiende. ¿Vos sentís que los jugadores son conscientes del riesgo de lesionarse o de la importancia de practicar correctamente los movimientos?

    \item[\textbf{Sebastián López:}] Hoy en día hay muchos entrenamientos funcionales para futbolistas, y eso también les abre la cabeza. Pero uno puede estar muy bien entrenado y muy bien preparado, y aun así el momento del partido es distinto. Hay otra presión y otra carga psicológica, y eso también tiene mucho que ver en el momento de la lesión. No por nada hoy en día, con los niveles de presión que se manejan en el fútbol argentino, hay tantas lesiones. Uno se pregunta: ``¿cómo se lesionó solo?''. Pero por ahí la cabeza juega un papel importante. A veces querés hacer un movimiento para el que no estabas preparado, querés adelantarte a la jugada y terminás haciendo un gesto que provoca la lesión.

    \item[\textbf{Entrevistador:}] Sí, es muy claro lo que comentás y súper entendible. ¿Te parecería útil que se pueda grabar el entrenamiento y luego recibir información sobre movimientos que, por ahí, son antiintuitivos o que podrían generar una lesión?

    \item[\textbf{Sebastián López:}] Sí, puede ser útil, siempre y cuando después haya un entrenamiento adaptado para mejorar esa carrera, ese giro o ese gesto de salto. Muchas veces son cuestiones técnicas las que tenés que mejorar. Uno de por sí no nace sabiendo correr correctamente. De chico corría para adelante, pero no sabía que tenía que apoyar de determinada manera, por ejemplo punta-talón. Eso tiene mucho que ver. Hoy en el fútbol profesional hay muchas herramientas que te dicen si tu marcha es buena o si tu pisada es buena. A la larga, eso repercute. Si no tenés una buena marcha o una buena pisada, puede repercutir en las articulaciones, y el problema puede aparecer más adelante.

    \item[\textbf{Entrevistador:}] Claro. Retomando esto que mencionás, ¿qué tendría que tener una herramienta así para que realmente la uses o le prestes atención? Por ejemplo, que te indique algún preparador qué aspecto se debe corregir.

    \item[\textbf{Sebastián López:}] Más allá de la herramienta, hoy en el fútbol profesional lo que tenés que buscar son ejercicios que generen ese cambio. Vos le podés decir a un jugador ``estás corriendo mal''. Se lo podés mostrar porque lo indica una máquina o un estudio, pero si le das esa información, después tenés que generar un entrenamiento que lo ayude a cambiar la marcha, mejorar técnicamente y sentir el cambio. Hay muchos jugadores que tienen problemas en los cuádriceps, que se les cargan, y muchas veces eso está encadenado con la postura o la cintura. Si tenés una mala postura, tal vez tenés que elongar más la espalda o la cintura. Entonces, se van haciendo ejercicios y el jugador empieza a notar mejoras. Hoy también existen herramientas como pilates, stretching o yoga, que a la larga ayudan y permiten sentir cambios. Entonces, creo que más que enfocarse solamente en encontrar el problema, aunque obviamente sin problema no hay ejercicio para hacer, lo importante es encontrar rápidamente el ejercicio que pueda ayudar a mejorar. En el nivel de competencia no tenés demasiado tiempo.

    \item[\textbf{Entrevistador:}] Buenísimo. Por lo que me comentás, todo depende de que, además de usar la herramienta, hay que prestar atención en las acciones posteriores que se deben realizar para que esto no ocurra. Clarísimo. Te agradezco mucho por todas las respuestas y por compartir tu experiencia. Con esto ya tengo información muy útil para el proyecto.

    \item[\textbf{Sebastián López:}] Dale, Ezequiel. Espero que te sirva para tu proyecto.
\end{description}

\anexo{Transcripción de la entrevista a Leonardo Lemos}
\begin{description}[leftmargin=0pt, itemindent=\parindent, labelsep=0.5em, itemsep=0pt, parsep=0pt, topsep=0pt]
    \item[\textbf{Entrevistador:}] Buenas, ¿cómo estás? Te agradezco por tomarte el tiempo. Empiezo contándote brevemente de qué trata el proyecto. La idea del proyecto es un sistema que utiliza análisis de video para observar movimientos del tren inferior en jugadores de fútbol, principalmente rodilla, tobillos y apoyos. La idea no es diagnosticar lesiones ni reemplazar al cuerpo médico, sino analizar videos de entrenamientos y competiciones para detectar posibles patrones de riesgo y generar información que pueda servir como apoyo para la prevención de lesiones. Para comenzar, ¿podrías hacer una breve presentación de quién sos y de tu trayectoria?

    \item[\textbf{Leonardo Lemos:}] Sí, sí. Mi nombre es Leonardo Lemos, soy entrenador de fútbol y en este momento estoy dirigiendo en el Club Los Andes. Trabajo como entrenador hace 19 años, desde los 30 años. Empecé en juveniles, donde estuve 10 años, y después ya llevo casi 10 años en el plantel de primera.

    \item[\textbf{Entrevistador:}] ¿Y en qué categorías y clubes trabajaste?

    \item[\textbf{Leonardo Lemos:}] Toda la historia de juveniles fue en Quilmes. Empecé con la novena, pasé por séptima, dirigí reserva, coordiné infantiles, coordiné todo el fútbol en general y después de ese proceso llegué a ser entrenador de primera. Después estuve en Sacachispas, Defensores de Belgrano, Colegiales, Brown de Madrín y ahora acá en Los Andes.

    \item[\textbf{Entrevistador:}] Gracias por el detalle, son muchísimos clubes. Y ahora que estás en Los Andes, ¿cómo es una semana normal de entrenamiento para el plantel?

    \item[\textbf{Leonardo Lemos:}] Tenemos el primer día pospartido, después del fin de semana, que puede ser lunes o martes según cuándo se juegue. Ese día se recuperan los jugadores que más minutos acumularon en el partido previo, y los demás realizan alguna actividad que les permita equilibrar cargas con respecto a ese grupo que jugó el fin de semana. La idea es tratar de que no queden desfasados en cuanto a la carga, más allá de que es muy difícil, porque obviamente la actividad propia del partido es irreemplazable. Pero se intenta que tengan una carga superior el primer día. El segundo día ya se junta todo el plantel otra vez. A mitad de semana empezamos a preparar el partido que viene y, sobre el cierre de la semana, trabajamos la parte de pelota parada.

    \item[\textbf{Entrevistador:}] Gracias por la explicación. ¿Se incluyen trabajos específicos de prevención de lesiones?

    \item[\textbf{Leonardo Lemos:}] Sí. Eso lo maneja el departamento médico en interacción con los profes, en las partes previas, en los preentrenamientos, antes del comienzo del entrenamiento en sí. Hay una coordinación entre ellos para realizar algunas actividades preventivas y, sobre todo, para ver casos puntuales de jugadores que necesitan algo específico. Se trata de abordar cada caso en base a sus necesidades. Después, obviamente, hay lineamientos generales en los preentrenamientos.

    \item[\textbf{Entrevistador:}] Claro. Entonces, por lo que comentás, esta parte de análisis y prevención está más vinculada a los preparadores físicos y al cuerpo médico que al cuerpo técnico directamente, ¿no?

    \item[\textbf{Leonardo Lemos:}] Sí. En realidad, esto se trabaja mediante la interacción entre el departamento médico y el departamento de preparación física. Entre ellos, según las respuestas y evaluaciones que van haciendo los profes con los GPS, los números y demás, van apareciendo datos. Esos datos se conectan con el departamento médico. También puede pasar que un jugador se acerque y manifieste alguna molestia, y ahí se va atacando cada caso de manera puntual, más allá de que hay lineamientos generales para todo el plantel.

    \item[\textbf{Entrevistador:}] Entiendo. ¿Vos considerás que los jugadores son conscientes del riesgo de lesiones, de los malos movimientos, malos apoyos, giros o frenadas?

    \item[\textbf{Leonardo Lemos:}] Sí. Es parte del juego. Está incluida la posibilidad de alguna molestia o alguna lesión. Por eso, de un tiempo a esta parte, se generó ese espacio de preentrenamiento, como para tratar de abordar un poco la prevención. Como te decía, también se ven cuestiones puntuales de cada jugador y cuestiones generales de los 30 que integran un plantel. Yo creo que el jugador cada vez es más profesional y, obviamente, si tiene alguna molestia o algo, interactúa con el kinesiólogo o con el doctor. A partir de eso salen ejercicios que después prepara el preparador físico, como para tratar de que llegue al entrenamiento con cierta prevención en cada situación.

    \item[\textbf{Entrevistador:}] Claro. También vi que fuiste jugador. ¿En qué año te retiraste?

    \item[\textbf{Leonardo Lemos:}] Me retiré en el 2007. Debuté en el 97 y me retiré en el 2007.

    \item[\textbf{Entrevistador:}] Desde tu retiro en 2007 hasta ahora, ¿ves un cambio en cuanto a la prevención de lesiones y lo que se analiza con los jugadores?

    \item[\textbf{Leonardo Lemos:}] Sí, en todos los sentidos se evoluciona constantemente. Hay mucha más información y mucha más posibilidad de obtener información. Entonces, desde ahí, todo puede ser más puntilloso y más específico. Quizás en mi época era más a ojo, y hoy hay datos más puntuales a partir de algunos artefactos que te permiten obtener información más fina y trabajar sobre eso. Así que sí, seguramente ha evolucionado un montón. Cada área también va mejorando la información que tiene, consiguiendo más información y trabajando con eso.

    \item[\textbf{Entrevistador:}] Claro, entiendo. Con el avance de la tecnología hay más posibilidades de medir y analizar. Desde tu experiencia, ¿qué importancia tiene poder detectar estos patrones antes de que aparezca una lesión?

    \item[\textbf{Leonardo Lemos:}] En realidad, es importante justamente para tratar de que no aparezca la lesión, en la medida en que se pueda. De hecho, como decía antes, los números de los GPS dan algunos datos con los que nosotros trabajamos. Quizás, si un jugador en un partido estuvo por sobre la media de lo que suele dar individualmente, o por sobre lo que se suele dar en esa posición, estamos atentos a ver si su recuperación necesita un día más, tal vez para no hacerlo trabajar en fatiga. Ese tipo de beneficios trae esto de la tecnología. Después hay lesiones que son fortuitas, pero hay cosas que se pueden ir viendo en base a los datos y, a partir de ahí, tratar de evitar que pase algo.

    \item[\textbf{Entrevistador:}] Seguro. ¿A vos te parecería útil que se pueda grabar un entrenamiento y recibir información sobre estos posibles riesgos? Entiendo que sería más para el preparador físico, pero me interesa tu mirada.

    \item[\textbf{Leonardo Lemos:}] Es que lo hacemos. Nosotros trabajamos también con algunas filmaciones, sobre todo en casos puntuales, por ejemplo cuando hay lesiones recurrentes. Se busca dónde está el error en la mecánica, quizás en la carrera. Entonces ahí aparecen las filmaciones en algún determinado momento de la actividad física, para ver si hay alguna situación que se pueda detectar en el video. Por ejemplo: ``acá está apoyando mal'' o ``acá está pisando distinto de lo que debería''. Eso sí lo hacemos, pero no lo hago yo directamente. Lo realiza más que nada la interacción entre el departamento médico y el departamento de preparación física. Sobre todo los kinesiólogos, que salen a la cancha y también preparan ejercicios preventivos o recuperatorios.

    \item[\textbf{Entrevistador:}] Gracias, me sirve tu punto de vista para ver también cómo trabajan en coordinación con el resto del club. Volviendo a la herramienta, ¿ves alguna barrera para utilizar estas herramientas? Por ejemplo, ¿los equipos de primera tienen alguna ventaja frente a equipos de otras divisiones?

    \item[\textbf{Leonardo Lemos:}] Puede ser lo económico. Quizás las posibilidades económicas, no solo en este aspecto sino en cualquiera, marcan alguna diferencia. La posibilidad de tener GPS hoy es bastante más amplia, por lo menos para las primeras dos categorías del fútbol argentino, quizás también para la tercera. En otra época eso era más lejano. Todo va avanzando, pero en definitiva lo económico termina siendo determinante en un montón de aspectos. No es más que eso. El que puede, seguramente va a estar interesado en tener la posibilidad de utilizar este tipo de herramientas porque ayudan.

    \item[\textbf{Entrevistador:}] Claro, entiendo. Para cerrar, ¿hay algo que no te haya preguntado y que te parezca importante sobre la prevención de lesiones, los entrenamientos o el uso de la tecnología en el fútbol?

    \item[\textbf{Leonardo Lemos:}] En mi caso, lo que puedo aportar es que es bienvenida. El aporte tecnológico es importante en cada aspecto. Obviamente, cada uno desde su función utiliza más algún dato o alguna información, y otra área u otro departamento utiliza otra. Lo importante es que, a partir de eso, podamos interactuar. Como te decía, por ahí los datos del GPS los interpreta el preparador físico y me los baja a mí de una manera lo más terrenal posible para que yo pueda entenderlos. A partir de ahí interactuamos: ``mirá, con este jugador hay que estar atento a esto''. Esa interacción es lo que se necesita para evitar, en la medida de lo posible, que pueda pasar algo. Obviamente es imposible evitarlo al 100\%. Lo mismo ocurre con el departamento médico. Tiene que haber interacción desde el cuerpo técnico hacia el departamento médico, tanto para la prevención como para la recuperación posterior en caso de que la lesión haya aparecido. En definitiva, lo importante es que la comunicación sea fluida.

    \item[\textbf{Entrevistador:}] Excelente, genial. Te agradezco mucho por todas las respuestas y por compartir tu experiencia.

    \item[\textbf{Leonardo Lemos:}] No, por favor. Estoy a disposición y que salga todo bien con el proyecto.
\end{description}

\anexo{Transcripción de la entrevista a Gonzalo Pereyra Metnyk}
\begin{description}[leftmargin=0pt, itemindent=\parindent, labelsep=0.5em, itemsep=0pt, parsep=0pt, topsep=0pt]
    \item[\textbf{Entrevistador:}] Bueno, a todos los entrevistados les menciono lo mismo, pero te agradezco por tomarte el tiempo y te cuento brevemente de qué trata el proyecto final. Estoy desarrollando una propuesta de sistema que utiliza inteligencia artificial y análisis de video para observar movimientos del tren inferior en jugadores de fútbol, principalmente rodilla, tobillo y apoyos. La idea no es diagnosticar lesiones ni reemplazar al kinesiólogo, al médico o al cuerpo técnico, sino analizar videos de entrenamientos y competiciones para detectar posibles patrones de riesgo y generar información que pueda servir como apoyo para la prevención de lesiones. La entrevista busca conocer tu experiencia profesional, cómo trabajás y, en particular, tu mirada sobre el uso de inteligencia artificial, machine learning y computer vision en este tipo de proyectos. También me interesa saber qué tan útil podría ser esta herramienta como apoyo en un contexto real. ¿Querés contarme quién sos, dónde trabajás y qué hacés en tu trabajo?

    \item[\textbf{Gonzalo Pereyra Metnyk:}] Sí, cómo no. Mi nombre es Gonzalo Pereyra Metnyk. Trabajo en inteligencia artificial hace aproximadamente ocho o nueve años. Arranqué trabajando en lo que antes también se asociaba a inteligencia artificial, que eran problemas de optimización, algoritmos genéticos y algoritmos de optimización en general. Después migré mucho a machine learning, específicamente deep learning. Trabajé mucho con computer vision, NLP y datos tabulares, dependiendo mucho del proyecto. Actualmente migré más hacia LLMs, que es lo que está más de moda en este momento. Trabajé en varios lugares. Entre los más importantes, creo que está Mercado Libre, donde estuve aproximadamente tres años. Ahí trabajé también con computer vision y bastante con NLP, sobre todo en cuestiones relacionadas con publicaciones de distintos objetos en el marketplace. También trabajé para una empresa de Holanda, con mucho foco en computer vision, cámaras y fotos de celulares. Había proyectos bastante interesantes. Actualmente estoy en una empresa importante vinculada a cuestiones impositivas, que tiene clientes muy grandes, como Google, Meta, Epic Games y otras empresas. Básicamente, el trabajo consiste en transformar data no estructurada en data estructurada. Eso incluye NLP para texto, computer vision para OCR, extracción de información de documentos de PDF escaneados, e incluso audio, porque a veces hay declaraciones en formato de audio. En definitiva, se trabaja con distintas fuentes para transformar esa información en datos que sirvan.

    \item[\textbf{Entrevistador:}] Gracias por la explicación. Me mencionaste que tenés conocimiento en machine learning y computer vision. Desde tu experiencia, ¿considerás viable aplicar estos conocimientos para analizar movimientos de futbolistas en video?

    \item[\textbf{Gonzalo Pereyra Metnyk:}] Sí. De hecho, el uso de deep learning en deporte no es algo nuevo. Hay muchos modelos que estiman poses en distintos deportes, como fútbol, básquet o yoga, entre otros. Existen muchos modelos open source que sirven para estimar poses. En deporte se usa bastante computer vision para sacar métricas de partidos, analizar jugadas, evaluar cómo fue un partido en general o incluso hacer predicciones según distintos factores. Así que sí, lo veo bastante viable y es algo que ya se está usando.

    \item[\textbf{Entrevistador:}] Perfecto. En relación con esto, ¿qué tecnologías o herramientas conocés para trabajar con estimación de postura corporal?

    \item[\textbf{Gonzalo Pereyra Metnyk:}] La estimación de postura corporal es algo bastante estudiado. Seguramente viste ejemplos donde aparece una persona moviéndose y el sistema marca distintos puntos en el cuerpo. Hay muchos modelos que estiman esos puntos. La idea es estimar cada punto y seguirlo a lo largo de todo el video. Yo te diría que arranques probando, porque en esta área depende mucho del caso y muchas veces hay que probar distintas alternativas hasta encontrar la que funciona. Literalmente, muchas veces se trabaja así. Me parece que YOLO puede ser un muy buen punto para arrancar. Es open source y es uno de los modelos más usados. Además, sirve bastante para cuestiones en tiempo real o videos que requieren inferencias rápidas. También hay modelos de Google para pose estimation. Uno de ellos, si no recuerdo mal, estaba muy asociado a videos de yoga. Otro modelo es MoveNet, que también es muy bueno y está optimizado para inferencias rápidas, por ejemplo en celulares. De todas formas, vuelvo a lo mismo: hay que probar y ver qué funciona mejor. Para un set de datos y una tarea específica, podés entrenar o probar muchos modelos, y no siempre se sabe de antemano cuál va a ser la mejor solución. Entonces, si probás YOLO y funciona bien, y después probás otro modelo más pesado que mejora apenas un 1\% la precisión, quizás no vale la pena cambiar. Una vez que encontrás un modelo que funciona razonablemente bien, conviene optimizar ese y avanzar.

    \item[\textbf{Entrevistador:}] Claro. Justo tocaste un tema que me interesa, que es la grabación con celular. ¿Qué tan confiable puede ser el análisis si el video se graba con una cámara común, con cierta resolución, o con un celular?

    \item[\textbf{Gonzalo Pereyra Metnyk:}] Depende. Ahí hay varias cuestiones. Primero, influye mucho la calidad del celular. No es lo mismo grabar a 30 FPS que a 60 FPS. En deporte, por la velocidad del movimiento, puede ser importante que el video capture bien los gestos. También depende mucho de cómo se use el celular. Si está quieto, en un trípode, filmando desde un lugar fijo, eso es más viable. Si hay una persona corriendo con el celular detrás del jugador, se complica muchísimo. En ese caso aparecen problemas de estabilización y calibración de cámara. Un punto muy importante es el set de entrenamiento. Cuando tengas un set de datos, los videos reales que uses deberían parecerse a los datos con los que entrenaste o validaste el modelo. Si entrenás con videos grabados con una cámara profesional en 4K, desde un ángulo ideal, después en la vida real el modelo va a esperar algo parecido. En cambio, si entrenás o validás con videos más parecidos a celulares comunes, con menor calidad, pocos FPS o condiciones menos ideales, con rayas o baja luminosidad, el modelo va a estar más preparado para ese contexto. Entonces, la confiabilidad depende mucho de con qué datos se entrena, con qué datos se valida y qué condiciones se replican después en el uso real.

    \item[\textbf{Entrevistador:}] Siguiendo con las limitaciones técnicas que pueden aparecer, ¿qué te parece más relevante? Por ejemplo, iluminación, ángulo de cámara, velocidad del movimiento o distancia.

    \item[\textbf{Gonzalo Pereyra Metnyk:}] Depende mucho del flujo de la aplicación. ¿Cómo te imaginás el funcionamiento?

    \item[\textbf{Entrevistador:}] La idea sería que un usuario cargue un video de un entrenamiento o competencia. Ese video se sube, se procesa y, aunque tarde un tiempo, no sería en tiempo real. Después, una vez procesado, el sistema marcaría puntos clave o movimientos relevantes para que un experto pueda validar si esos movimientos podrían representar algún riesgo.

    \item[\textbf{Gonzalo Pereyra Metnyk:}] Eso aligera bastante el problema, porque si no es en tiempo real puede tardar más en procesar. Respecto a las limitaciones como iluminación, ángulo o distancia, en la vida real se solucionan principalmente entrenando o validando con más variedad de datos. Si yo tengo solo fotos tuyas de frente, el modelo va a aprender muy bien esa situación. Pero si después en la vida real te filman de costado, desde atrás o con poca luz, el modelo necesita haber visto ejemplos similares durante el entrenamiento o la validación. El problema es que entrenar con muchos datos, especialmente imágenes o videos, puede ser pesado. Se suele necesitar GPU y no siempre es sencillo. Por eso, generalmente se usa un modelo ya entrenado y se aplica transfer learning o fine tuning. Es decir, se toma un modelo que ya aprendió características generales, por ejemplo detectar personas o partes del cuerpo, y se lo adapta a tu caso específico. En modelos de imágenes, las primeras capas suelen detectar formas más generales, como bordes o contornos, y las últimas capas se especializan en características más específicas. Entonces no necesitás entrenar todo desde cero. Conviene partir de un modelo que ya entiende qué es una persona o una postura, y ajustarlo a tu caso: jugadores de fútbol, movimientos específicos, posibles lesiones o patrones de riesgo. Mi recomendación sería separar el problema en etapas. Por ejemplo, primero un modelo que estime la pose y entienda qué está haciendo el jugador. Después, otro modelo o componente que, en base a esa pose, analice posibles problemas o patrones de riesgo. Si hacés que un solo modelo haga todo, podés ganar performance porque corrés un único modelo, pero quizás perdés precisión o explicabilidad. Si separás en varios modelos o etapas, podés ganar precisión, pero aumenta la complejidad. Siempre hay un trade-off. En definitiva, hay que probar con tus datos y tu problema concreto. No se puede decir de antemano cuál es la mejor solución.

    \item[\textbf{Entrevistador:}] Entiendo. Para un primer prototipo, ¿conviene analizar ejercicios controlados antes que jugadas reales?

    \item[\textbf{Gonzalo Pereyra Metnyk:}] Si en el uso real vas a analizar ejercicios controlados, entonces usá datos controlados. Si en la realidad vas a analizar jugadas reales, entonces el set de validación tiene que incluir jugadas reales. Yo me enfocaría primero en la validación. Tu set de evaluación tiene que ser lo más parecido posible a lo que vas a usar en la vida real. Si en la vida real son jugadas reales, entonces el set de validación tiene que tener jugadas reales sí o sí. Después podés entrenar con datos reales, y si te sirve, agregar jugadas controladas. También podés agregar más variedad o incluso datos de deportes parecidos, si aportan información útil. Pero lo importante es validar siempre contra datos representativos de la realidad. La vida real cambia mucho. Las imágenes no suelen ser perfectas. En un dataset capaz tenés una pose muy clara, pero en la realidad el jugador está girando, cayéndose, cambiando de dirección o haciendo movimientos que el modelo nunca vio. Por eso, el set de validación es clave.

    \item[\textbf{Entrevistador:}] ¿Te parece correcto plantear el sistema como detección de patrones de riesgo y no como predicción o diagnóstico de lesiones? Te comento que considero como diferencia de estos. En un caso, la idea sería mostrar que ciertos patrones que está haciendo el jugador podrían generar riesgo de lesión. En cambio, una predicción o diagnóstico sería decir, por ejemplo, que hay un porcentaje de probabilidad de que tenga una rotura de ligamento en los próximos seis meses.

    \item[\textbf{Gonzalo Pereyra Metnyk:}] Entiendo. Yo creo que las dos cosas van bastante de la mano. Para detectar un patrón de riesgo, de alguna manera tenés que relacionar ese movimiento con la posibilidad de que genere un problema. Podrías hacer ambas cosas tranquilamente. De hecho, algo interesante sería predecir si el jugador podría tener una lesión y, si la confianza del modelo es alta, analizar los movimientos que hizo para asociar ciertos gestos con mayor probabilidad de lesión. De todas formas, entiendo que para un proyecto académico puede ser más prudente hablar de patrones de riesgo, sobre todo si no se busca hacer un diagnóstico médico.

    \item[\textbf{Entrevistador:}] Claro. ¿Qué riesgos éticos o de privacidad considerás que habría que tener en cuenta al trabajar con videos de deportistas o personas en general?

    \item[\textbf{Gonzalo Pereyra Metnyk:}] El tema de los datos es muy complejo. En la práctica, muchas empresas analizan una gran cantidad de información. Por eso, es importante que los datos estén protegidos y que no se filtren fuera de la organización. En este caso, una medida importante sería anonimizar los datos. Muchas veces no hace falta saber el nombre del jugador; alcanza con analizar el movimiento de una persona realizando determinada acción. Si se anonimiza la información, es más fácil reducir el riesgo de que afecte al jugador.

    \item[\textbf{Entrevistador:}] Pero en cuanto al jugador específicamente, ¿no considerás que le puede generar problemas si se filtra la información, que hay clubes interesados en que se revele esa información por más que sea propia del club?

    \item[\textbf{Gonzalo Pereyra Metnyk:}] Sí, sí. También puede haber riesgos específicos en el deporte. Por ejemplo, si se filtra información sobre una posible lesión o un riesgo físico, eso podría afectar la carrera de un futbolista, una transferencia o la estrategia de otro equipo. Entonces, la información debería quedar dentro del club o del equipo autorizado. El problema no es solo ético, también es de seguridad de la información. Si se filtran datos de una empresa, de un banco o de un club, puede haber consecuencias importantes. Por eso, habría que cuidar mucho el acceso, la privacidad, la anonimización y el almacenamiento de los datos.

    \item[\textbf{Entrevistador:}] Claro, tiene sentido. Para cerrar, ¿hay alguna recomendación técnica o metodológica que consideres importante para este proyecto?

    \item[\textbf{Gonzalo Pereyra Metnyk:}] Sí, varias. Primero, no pensaría el sistema como algo exclusivamente de machine learning o exclusivamente de reglas. Podés combinar varias fuentes de información. Un modelo de machine learning procesa señales. Tenés una entrada, que son los datos, un modelo que procesa esa información y una salida. Pero puede pasar que un modelo funcione bien sin usar demasiado una variable que para un kinesiólogo o un especialista es muy importante, como el ángulo de la rodilla. Entonces, si un kinesiólogo te dice que cierto ángulo o cierto movimiento es relevante, esa información también se puede incorporar. Podés combinar lo que dice el modelo con reglas, parámetros biomecánicos, ángulos, información temporal u otros datos. No son excluyentes uno de otros. Por ejemplo, no es lo mismo analizar un movimiento al principio del entrenamiento que al final, cuando el jugador ya está fatigado. El tiempo o el nivel de fatiga podrían ser variables importantes. No tiene que ser todo blanco o negro. Podés armar una función o un sistema que combine distintas señales. La recomendación principal sería probar todo lo que se te ocurra, siempre validando correctamente. Mientras el set de validación esté bien construido y no haya filtración de datos entre entrenamiento y validación, podés experimentar con distintas combinaciones. Otra recomendación es limitar mucho el alcance. Me parece bien que te enfoques en tren inferior y futbolistas. No te compliques intentando resolver todo. No entrenes un modelo de pose desde cero, porque no es sencillo y puede ser muy costoso. Enfocate en tu área específica, aprovechá modelos existentes y validá bien tu solución.

    \item[\textbf{Entrevistador:}] Buenísimo. Con esto terminaría la entrevista. Te agradezco mucho por haberte tomado el tiempo y por todas las recomendaciones. Me sirve mucho para el proyecto.

    \item[\textbf{Gonzalo Pereyra Metnyk:}] Un lujo. Muy buena tu idea y éxitos con el proyecto.
\end{description}